\section{Discussion and Limitations}

Baihu provides a large-scale and standardized data foundation for embodied intelligence research. Its main strengths lie in its scale, fully human-teleoperated collection protocol, multi-platform coverage, preservation of platform-specific observation and action schemas, and conversion into the LeRobot v2.1 format. With \textbf{9521.75 hours}, \textbf{513575 episodes}, and more than \textbf{1.03 billion frames} of robot data, Baihu supports pretraining regimes that are difficult to study using smaller task-specific datasets.

The offline zero-shot benchmark provides initial evidence that Baihu is useful for embodied foundation model pretraining. The evaluation datasets are external to Baihu and are excluded from Baihu training. Across 13 platforms and 42 paired task-dataset records, the Baihu-trained GR00T N1.6 checkpoint substantially reduces both Normalized ALL MSE and Normalized Joint MSE compared with the no-training checkpoint. Each task-level result is averaged over five evaluation trajectories, and both checkpoints are evaluated on the same trajectories. These results suggest that Baihu pretraining transfers effectively to separately collected robot demonstration data.

At the same time, the current benchmark should be interpreted as an offline action-prediction evaluation rather than a complete measure of real-world robot performance. Open-loop normalized MSE measures agreement with reference teleoperation actions, but it does not capture closed-loop recovery, compounding errors, contact-rich interaction, or task success under real robot execution. Moreover, the benchmark demonstrates dataset-level transfer because the evaluation datasets are external to Baihu, but it should not be interpreted as a held-out-platform experiment when the evaluated robot platform is also represented in Baihu.

Another limitation is that the current paper focuses on a single primary benchmark setting: a GR00T N1.6 checkpoint comparison under offline open-loop evaluation. Additional baselines such as ACT, Diffusion Policy, or other vision-language-action policies would make the benchmark more comprehensive. Data-scaling experiments, held-out-platform evaluation, and closed-loop robot rollouts would further clarify which aspects of Baihu contribute most to model performance. The current benchmark also uses an unweighted mean across task-dataset records; future work could additionally report trajectory-weighted results and uncertainty estimates across evaluation trajectories.